\documentclass[11pt]{article}
\usepackage[margin=1in]{geometry}
\usepackage{amsmath,amssymb,amsthm,booktabs,longtable,array}
\usepackage[T1]{fontenc}
\usepackage{lmodern,microtype}
\usepackage[colorlinks=true,linkcolor=blue,citecolor=blue,urlcolor=blue]{hyperref}
\pdfinfoomitdate=1
\pdftrailerid{}
\newtheorem{theorem}{Theorem}
\newtheorem{lemma}[theorem]{Lemma}
\newtheorem{proposition}[theorem]{Proposition}
\newtheorem{corollary}[theorem]{Corollary}
\theoremstyle{definition}\newtheorem{definition}[theorem]{Definition}
\newcommand{\R}{\mathbb R}
\newcommand{\C}{\mathbb C}
\newcommand{\E}{\mathbb E}
\newcommand{\He}{\operatorname{He}}
\newcommand{\sgn}{\operatorname{sgn}}
\newcommand{\erf}{\operatorname{erf}}
\newcommand{\erfc}{\operatorname{erfc}}
\newcommand{\dd}{\,d}
\newcommand{\norm}[1]{\left\lVert#1\right\rVert}
\newcommand{\abs}[1]{\left\lvert#1\right\rvert}
\newcommand{\fall}[2]{(#1)_{#2}}
\title{A certified upper bound for the real Grothendieck constant}
\author{}
\date{}
\begin{document}
\maketitle
\begin{abstract}
An explicit Gaussian rounding scheme gives
\[
 K_G^{\R}\leq \frac{\pi}{2(0.8818276493988)}
 <1.781296297373.
\]
The scheme uses a degree-nine polynomial threshold and a degree-eighty-one
tensor preprocessing polynomial. The proof combines a coefficient inequality
valid for every finite real matrix, interval enclosures of all 63252 Gaussian
moments required through degree 501, and a bound on the fourth angular
derivative on the entire unit circle. We give the Gaussian identities,
quadrature remainders, boundary convergence argument, and exact rational
comparisons underlying the certificate. All polynomial coefficients and the
finite data needed to reproduce the calculation accompany the proof.
\end{abstract}

\section{Statement and normalization}
For a finite real matrix $M=(M_{ij})_{1\leq i\leq m,1\leq j\leq n}$, put
\begin{align*}
 S(M)&=\max_{\varepsilon_i,\delta_j\in\{-1,1\}}
       \abs{\sum_{i,j}M_{ij}\varepsilon_i\delta_j},\\
 V(M)&=\max_{\norm{u_i}=\norm{v_j}=1}
       \abs{\sum_{i,j}M_{ij}\langle u_i,v_j\rangle}.
\end{align*}
The vectors are real. Their span has dimension at most $m+n$, so the latter
maximum can be taken in $\R^{m+n}$ and is attained. The real Grothendieck
constant is the least $K$ for which $V(M)\leq K S(M)$ for every such matrix.

\begin{theorem}\label{thm:main}
For every pair of positive integers $m,n$ and every $M\in\R^{m\times n}$,
\[
 V(M)\leq\frac{\pi}{2(0.8818276493988)}S(M).
\]
In particular, $K_G^{\R}<1.781296297373<1.78130$.
\end{theorem}

The normalization throughout is $H=(\pi/2)\E[fg]$ for real Gaussian
rounding. Gaussian rounding and tensor preprocessing originate in work of
Krivine~\cite{krivine}; multivariate rounding was developed further by
Braverman, Makarychev, Makarychev, and Naor~\cite{bmmn}. The combination of
coordinate-dependent preprocessing, Hermite coefficients, a coefficient
margin, and a Sobolev estimate follows methods used by Saha et
al.~\cite{saha}. The argument below proves the needed statements for the
explicit polynomials used here; it does not use an upper-bound theorem from
those papers as a premise.

For comparison, the numerical inequalities in Proposition 6.3 of
\cite{saha}, version 2, imply the upper bound
\begin{equation}\label{eq:comparison-saha}
 \frac{\pi}{2\left(0.881573822049-0.000011328860
              -14.443/\sqrt{10\cdot251^5}\right)}
 =1.78184132423347648\ldots.
\end{equation}
Li et al.~\cite{li}, version 3, separately report $1.7813319810625639$
as an internally tested upper bound, expressly stating that the authors had
not verified those certificates and that the reported claims were not
theorems. The bound in Theorem~\ref{thm:main} is numerically smaller than
both quantities. No historical priority or exact-value assertion is made.

\section{A coefficient inequality for every finite matrix}
Let $\gamma$ be standard Gaussian measure on $\R$, with density
$\phi(x)=(2\pi)^{-1/2}e^{-x^2/2}$. Define the probabilists' Hermite
polynomials and their normalization by
\[
 e^{tx-t^2/2}=\sum_{k\geq0}\He_k(x)\frac{t^k}{k!},
 \qquad h_k=\frac{\He_k}{\sqrt{k!}}.
\]
Taking Gaussian expectations of two generating functions proves
$\int h_jh_k\dd\gamma=\mathbf1_{j=k}$. These polynomials are complete
in $L^2(\gamma)$: if $u$ is orthogonal to all polynomials, then
$\int u(x)e^{zx}\dd\gamma(x)$ is entire, by Cauchy--Schwarz uniformly on
compact subsets of $\C$, and all its derivatives at zero vanish.
Its values on the imaginary axis are the Fourier transform of the finite
measure $u\dd\gamma$, so Fourier uniqueness gives $u=0$. The same argument,
or iteration in each variable, gives the tensor product basis in two
dimensions. If $(X,Y)$ are standard real Gaussians with covariance $t$, the
generating functions give
\begin{equation}\label{eq:hermite-correlation}
 \E[h_j(X)h_k(Y)]=\mathbf1_{j=k}t^j,\qquad -1\leq t\leq1.
\end{equation}
The endpoints follow either directly from $Y=\pm X$ or by continuity.

Let $L$ be the following odd polynomial:
\begin{equation}\label{eq:L}
 L(x)=\frac{3188129h_1(x)+110767301h_3(x)+13189406h_5(x)
                   -5120076h_7(x)-1681617h_9(x)}{10^9}.
\end{equation}
Set $f(w,x)=\sgn(w+L(x))$ and $g(w,x)=\sgn(w-L(x))$.
The value of the sign at zero is immaterial, since conditioning on $X$
shows that each threshold has Gaussian measure zero. Write
\[
 A_{jk}=\int_{\R^2}f(w,x)h_j(w)h_k(x)\dd\gamma(w)\dd\gamma(x).
\]
Since $f^2=1$, Parseval gives $\sum_{j,k}A_{jk}^2=1$. Oddness under
$(w,x)\mapsto(-w,-x)$ gives $A_{jk}=0$ when $j+k$ is even.
Moreover, $g(w,x)=f(w,-x)$, so its corresponding coefficient is
$(-1)^k A_{jk}$.

For independent coordinate pairs $(W,W')$ and $(X,X')$ with covariances
$\alpha$ and $\beta$, respectively, (\ref{eq:hermite-correlation}) gives
\begin{equation}\label{eq:F}
 F(\alpha,\beta):=\frac\pi2\E[f(W,X)g(W',X')]
 =\frac\pi2\sum_{j,k\geq0}(-1)^k A_{jk}^2\alpha^j\beta^k.
\end{equation}
This identity follows first by finite Hermite approximation and the
$L^2$ contraction property of conditional expectation. The series is
absolutely and uniformly convergent on the closed complex bidisk, since
$\sum A_{jk}^2=1$. Thus it also defines a continuous function there,
holomorphic in the open bidisk.

Let $P(t)=\sum_{d=1,3,\ldots,81}p_dt^d$ be the exact rational polynomial
in Appendix~\ref{app:polynomial}. Its coefficient norm is
\begin{equation}\label{eq:Pnorm}
 \rho_0:=\sum_d|p_d|
 =\frac{499995394004774337}{500000000000000000}<1.
\end{equation}
Define
\begin{equation}\label{eq:H}
 H(t)=F(P(t),t)=\sum_{n\geq0}b_nt^n.
\end{equation}
The coefficients are real and vanish for even $n$. In the Banach algebra
of absolutely summable power series, substitution in (\ref{eq:F}) gives
\begin{equation}\label{eq:Habsolute}
 \sum_n|b_n|\leq\frac\pi2\sum_{j,k}A_{jk}^2\norm{P}_1^j
 \leq\frac\pi2.
\end{equation}
This also justifies substitution and evaluation at $t=\pm1$.

\begin{lemma}\label{lem:rounding}
If the coefficients in (\ref{eq:H}) satisfy
$\Gamma=b_1-\sum_{n\geq3,\ n\ {\rm odd}}|b_n|>0$, then
$V(M)\leq\pi S(M)/(2\Gamma)$ for every finite real matrix $M$.
\end{lemma}
\begin{proof}
For real unit vectors $u,v$, form the direct sums
\[
 U_P(u)=\bigoplus_d\sqrt{|p_d|}\,u^{\otimes d}
               \ \oplus\sqrt{1-\rho_0}\,e_L,
 \quad
 V_P(v)=\bigoplus_d\sgn(p_d)\sqrt{|p_d|}\,v^{\otimes d}
               \ \oplus\sqrt{1-\rho_0}\,e_R,
\]
where $e_L,e_R$ are orthogonal unit vectors in extra coordinates.
Both vectors have norm one and their cross inner product is
$P(\langle u,v\rangle)$. This is a finite-dimensional real construction;
negative polynomial coefficients cause no loss of norm or positivity of
the joint Gaussian covariance. Use a standard Gaussian vector in this
space to obtain all $W_i,W'_j$ simultaneously. In an independent Gaussian
space use the original $u_i,v_j$ to obtain $X_i,X'_j$. Each endpoint has
independent standard $W$ and $X$, and each cross pair has correlations
$P(t_{ij})$ and $t_{ij}$, where $t_{ij}=\langle u_i,v_j\rangle$.

Choose vectors attaining $V(M)$, reversing all left vectors if needed so
that $\sum M_{ij}t_{ij}=V(M)$. For every positive integer $n$, the unit
vectors $u_i^{\otimes n},v_j^{\otimes n}$ show that
$|\sum M_{ij}t_{ij}^n|\leq V(M)$. Absolute summability in
(\ref{eq:Habsolute}) and finiteness of $M$ therefore imply
\[
 \E\sum_{i,j}M_{ij}f(W_i,X_i)g(W'_j,X'_j)
 =\frac2\pi\sum_{i,j}M_{ij}H(t_{ij})
 \geq\frac2\pi\Gamma V(M).
\]
For each Gaussian outcome the signed sum is at most $S(M)$. Taking
expectations proves the assertion, also when $V(M)=0$.
\end{proof}

\section{The complete coefficient calculation through degree 501}
\subsection{One-dimensional moments and exact finite algebra}
Differentiating the Hermite generating function gives
\begin{equation}\label{eq:hermite-recurrence}
 h_{k+1}(x)=\frac{xh_k(x)-\sqrt{k}h_{k-1}(x)}{\sqrt{k+1}},
 \qquad h'_k(x)=\sqrt{k}\,h_{k-1}(x).
\end{equation}
For $a\geq0$, put $q_a(u)=\int\sgn(w+u)h_a(w)\dd\gamma(w)$.
The identity $(h_{a-1}\phi)'=-\sqrt a\,h_a\phi$, including the
vanishing boundary term at infinity, gives
\begin{align}
 q_0(u)&=\erf(u/\sqrt2),\nonumber\\
 q_a(u)&=\sqrt{\frac2\pi}\,e^{-u^2/2}
                      \frac{h_{a-1}(-u)}{\sqrt a}\quad(a\geq1),\label{eq:q}\\
 q_{a+1}(u)&=\frac{-u\sqrt a\,q_a(u)-(a-1)q_{a-1}(u)}{\sqrt{a(a+1)}}
                       \quad(a\geq1).\nonumber
\end{align}
Consequently
\begin{equation}\label{eq:Aintegral}
 A_{ab}=\int_\R h_b(x)q_a(L(x))\dd\gamma(x).
\end{equation}
For real $u$, Cauchy--Schwarz gives $|q_a(u)|\leq1$ for every $a$.
Thus the entire omitted real Gaussian tail satisfies
\begin{equation}\label{eq:moment-tail}
 \abs{\int_{|x|>16}h_b(x)q_a(L(x))\dd\gamma(x)}
 \leq\sqrt{\erfc(16/\sqrt2)}.
\end{equation}
For $a+b$ odd the integrand in (\ref{eq:Aintegral}) is even.

We calculate every moment in
\[
 \mathcal I=\{(a,b)\in\mathbb Z_{\geq0}^2:a+b\leq501,
                                      \ a+b\text{ odd}\}.
\]
There are $\sum_{r=0}^{250}(2r+2)=63252$ such moments. Since $P(0)=0$,
$P(t)^at^b$ has order at least $a+b$. It follows rigorously that no
moment outside $\mathcal I$ contributes to any coefficient of degree at
most 501. In particular,
\begin{equation}\label{eq:head-formula}
 b_n=\frac\pi2\sum_{\substack{a,b\geq0\\a+b\leq n}}
                (-1)^b A_{ab}^2[t^{n-b}]P(t)^a
       \quad(1\leq n\leq501).
\end{equation}
All powers of $P$ may be computed in $\mathbb Q[t]/(t^{502})$; reduction
modulo $t^{502}$ is a ring homomorphism. There are 251 odd coefficients
$b_1,b_3,\ldots,b_{501}$, including $b_1$.

\subsection{Nodes, weights, and a uniform quadrature remainder}
We use 128 equal panels on $[0,16]$, with half-width $h=1/16$, and
112-point Gauss--Legendre quadrature on each panel. To specify why the
quadrature is rigorous, let $\mathcal P_s$ denote the Legendre polynomial
with $\mathcal P_s(1)=1$. Rodrigues' formula, followed by integration by
parts, gives orthogonality to polynomials of degree less than $s$ and
$\int_{-1}^1\mathcal P_s^2=2/(2s+1)$. There are 112 disjoint rational
brackets inside $(-1,1)$, with exactly checked opposite signs of
$\mathcal P_{112}$ at their endpoints. Degree 112 therefore implies
that they contain all its roots, each simple. For a root $z_i$, the
Lagrange polynomial
$\ell_i(x)=\mathcal P_{112}(x)/((x-z_i)\mathcal P'_{112}(z_i))$
gives weights $w_i=\int_{-1}^1\ell_i^2>0$. Polynomial division by
$\mathcal P_{112}$ shows that $\ell_i^2-\ell_i$ has integral zero,
so $w_i=\int_{-1}^1\ell_i$. For any polynomial of degree at most 223,
division gives a multiple of $\mathcal P_{112}$ with quotient of degree
at most 111, plus an interpolating remainder of degree at most 111.
Orthogonality and Lagrange interpolation therefore prove exactness
through degree 223. The
Christoffel--Darboux identity, obtained by telescoping the three-term
Legendre recurrence, gives
\begin{equation}\label{eq:weights}
 w_i=\frac{2}{(1-z_i^2)\mathcal P'_{112}(z_i)^2},\qquad
 \sum_iw_i=2.
\end{equation}
For completeness, the diagonal identity used here is
\[
 \sum_{k=0}^{s-1}\frac{2k+1}{2}\mathcal P_k(x)\mathcal P_k(y)
 =\frac{s}{2}\frac{\mathcal P_s(x)\mathcal P_{s-1}(y)
                         -\mathcal P_{s-1}(x)\mathcal P_s(y)}{x-y};
\]
at a root of $\mathcal P_s$, use
$(1-z_i^2)\mathcal P'_s(z_i)=s\mathcal P_{s-1}(z_i)$.

All root and weight evaluations use enclosing intervals. A midpoint
derivative enclosure is enlarged by
$\sup_{[-1,1]}|\mathcal P''_{112}|$ times the bracket radius. The exact
bound is
\[
 \sup_{[-1,1]}|\mathcal P_s''|
 \leq\frac{(s-1)s(s+1)(s+2)}8.
\]
Indeed $\mathcal P_s''=3C_{s-2}^{5/2}$; the generating function
$(1-2xt+t^2)^{-5/2}$ and its factorization for $x=\cos\theta$ bound
the absolute value of every coefficient by its value at $x=1$.
Exact rational sign checks, disjointness checks, and further rational
bisections independently enclose all nodes and the weights
(\ref{eq:weights}). Initial floating-point guesses play no evidentiary role.

We next bound the integrand on full complex ellipses, uniformly for all
moments. For $0<\tau<1$ and $z=x+iy$, Mehler's identity implies
\begin{equation}\label{eq:mehler}
 \sum_{k\geq0}\tau^k|h_k(z)|^2
 =(1-\tau^2)^{-1/2}
     \exp\left(\frac{\tau x^2}{1+\tau}
                    +\frac{\tau y^2}{1-\tau}\right).
\end{equation}
One direct justification, valid for complex $z$, uses the entire kernel
\[
 K_s(z,u)=\frac1{\sqrt{1-s^2}}
 \exp\left(\frac{2szu-s^2(z^2+u^2)}{2(1-s^2)}\right),\quad s=\sqrt\tau.
\]
Gaussian integration of its product with the Hermite generating function
shows that its Hermite coefficients in $u$ are $s^kh_k(z)$.
It belongs to $L^2(\gamma(du))$ for each complex $z$. Completing the
square in $\int|K_s(z,u)|^2\dd\gamma(u)$ and using Parseval gives
(\ref{eq:mehler}), including its convergence assertion.

Fix $N=501$, $\tau=99/100$, and $\rho=3/2$, and let
$F_N=\tau^{-N/2}(1-\tau^2)^{-1/4}$. Dropping a term from
(\ref{eq:mehler}) gives, for $b\leq N$,
\begin{equation}\label{eq:hphi-bound}
 |h_b(z)\phi(z)|\leq\frac{F_N}{\sqrt{2\pi}}
 \exp\left(-\frac{(\Re z)^2}{2(1+\tau)}
                  +\frac{(\Im z)^2}{2(1-\tau)}\right).
\end{equation}
Equation (\ref{eq:q}) gives the same exponential bound for $q_a(u)$,
with prefactor $\sqrt{2/\pi}F_N$, when $1\leq a\leq N$.
For $q_0$, integrate its derivative along a vertical segment to obtain
\begin{equation}\label{eq:qzero-bound}
 |q_0(u)|\leq1+\sqrt{2/\pi}|\Im u|
                   \exp\bigl(((\Im u)^2-(\Re u)^2)/2\bigr).
\end{equation}

For a panel centered at $c=(2j+1)h$, define
\[
 x_r=h(\rho+\rho^{-1})/2,\quad y_r=h(\rho-\rho^{-1})/2,
 \quad R=(x_r^2+y_r^2)^{1/2}.
\]
The closed Bernstein ellipse for this panel lies in
$|\Re z-c|\leq x_r$, $|\Im z|\leq y_r$, $|z-c|\leq R$.
Exact polynomial Taylor expansion bounds $|\Re L(z)|$ below by
\[
 U_c=\max\left\{0,|L(c)|-
                       \sum_{k=1}^9\frac{|L^{(k)}(c)|R^k}{k!}\right\},
\]
and bounds $|\Im L(z)|$ above by
\[
 V_c=y_r\sum_{k=0}^8\frac{|L^{(k+1)}(c)|R^k}{k!}.
\]
The latter bound follows by integrating $L'$ along the vertical segment
from $\Re z$ to $z$; $L$ is real on the real axis. Also
$|\Re z|\geq X_c:=\max\{0,c-x_r\}$. Thus a uniform upper bound on
$|h_b(z)q_a(L(z))\phi(z)|$ on the entire ellipse is
\begin{align}\label{eq:Mj}
 M_c={}&\frac{F_N}{\sqrt{2\pi}}
 e^{-X_c^2/(2(1+\tau))+y_r^2/(2(1-\tau))}\nonumber\\
 &\quad\cdot\max\left\{\sqrt{2/\pi}F_N
 e^{-U_c^2/(2(1+\tau))+V_c^2/(2(1-\tau))},
 1+\sqrt{2/\pi}V_c e^{(V_c^2-U_c^2)/2}\right\}.
\end{align}

If an analytic function has absolute value at most $M$ on a Bernstein
ellipse of parameter $\rho$, its Chebyshev coefficients of positive
degree have absolute value at most $2M\rho^{-k}$, by the Cauchy formula
after $z=(w+w^{-1})/2$. Truncation after degree 223 therefore has uniform
error at most $2M\rho^{-224}/(1-\rho^{-1})$ on $[-1,1]$.
Both integration and the positive Gauss rule have functional norm two
on that interval. Exactness through degree 223 now bounds their
difference by $8M\rho^{-224}/(1-\rho^{-1})$. Scaling each panel and
doubling the half-line integral gives the uniform complete error
\begin{equation}\label{eq:uniform-error}
 E_A=\frac{16h\rho^{-224}}{1-\rho^{-1}}
           \sum_{j=0}^{127}M_{(2j+1)h}
           +\sqrt{\erfc(16/\sqrt2)}<4.121\cdot10^{-22}.
\end{equation}
The terms are evaluated with outward rounding, including all square
roots and transcendental functions. The first term is less than
$4.1200500\cdot10^{-22}$ and the second less than
$3.574565\cdot10^{-29}$. This bound uses $N=501$ everywhere in
$F_N$; it is not the degree-251 quadrature error.

\subsection{Finite certificate and resulting coefficient inequalities}
The accompanying certificate contains the 112 node enclosures, 128
ellipse bounds, and 32 blocks of four panels each. Each block contains
every index in $\mathcal I$. The blocks are summed with interval
arithmetic, after which the symmetric error interval $[-E_A,E_A]$ is
added once to each moment. No panel is omitted or repeated. The moment
calculation uses 384-bit working precision, with 640 bits for initial
node refinement. All decimal intervals are serialized as
$(m\pm r)10^e$ with integer $m,r,e$ and are read with outward rounding.

Formula (\ref{eq:head-formula}) is then evaluated in two ways: by
interval power-series arithmetic and by exact rational powers of $P$
followed by interval evaluation of the moment factors. Both contain
every odd coefficient through 501 with its factor $(-1)^b$ unchanged.
The exact rational implementation gives the strict outward inequalities
\begin{equation}\label{eq:head-bounds}
 b_1>B:=0.88182999999999823752,
 \qquad \sum_{\substack{3\leq n\leq501\\n\ {\rm odd}}}|b_n|
       <h_*:=0.000000001054.
\end{equation}
For orientation, the interval sum before this rational enlargement is
$[1.0537295443\cdot10^{-9}\ \pm\ 5.65\cdot10^{-20}]$.
Only the outward inequalities (\ref{eq:head-bounds}) are used below.

\section{Regularity and a bound on the entire unit circle}
Let $D=z\,d/dz$. Since $d/d\theta=iD$ on a circle, the fourth
angular derivative is exactly $D^4$. We will prove
\begin{equation}\label{eq:C4}
 \left(\frac1{2\pi}\int_0^{2\pi}|D^4H(e^{i\theta})|^2\dd\theta\right)^{1/2}
 <C_*:=24744.587741.
\end{equation}
Two different estimates for the mixed derivatives of $F$ are used on
closed angular intervals. Before giving the numerical estimates, we
establish the boundary regularity needed to interpret this integral and
to pass from it to coefficient bounds.

\subsection{Gaussian smoothing in one coordinate}
For $0<r<1$, put $v=1-r$ and
\[
 G_r(w,x)=\E_Z f(\sqrt r\,w+\sqrt{1-r}\,Z,x)
 =\erf\left(\frac{\sqrt r\,w+L(x)}{\sqrt{2v}}\right).
\]
Here $Z$ is an independent standard real Gaussian. The generating
function or (\ref{eq:hermite-correlation}) shows that $G_r$ has Hermite
coefficient $r^{j/2}A_{jk}$. For nonnegative $p,q$ with $p+q\geq1$ define
\begin{equation}\label{eq:Spq}
 S_{pq}(r)=r^{-p}\norm{\partial_w^p\partial_x^qG_r}_{L^2(\gamma^2)}^2
 =\sum_{j,k}\fall{j}{p}\fall{k}{q}r^{j-p}A_{jk}^2,
\end{equation}
where $(j)_p=j!/(j-p)!$ if $j\geq p$ and is zero otherwise.
To justify this identity without assuming differentiability of $f$, note
that every fixed derivative of the explicit smooth function $G_r$ is a
polynomial in $w,x$ times a Gaussian in $\sqrt r\,w+L(x)$. It and all
lower derivatives have at most polynomial growth and belong to
$L^2(\gamma^2)$. Gaussian integration by parts therefore has vanishing
boundary terms. It identifies the Hermite coefficients of the derivative
as the shifted coefficients of $G_r$ multiplied by
$\sqrt{(j)_p(k)_q}$. Completeness and Parseval prove
(\ref{eq:Spq}), including finiteness.

Termwise differentiation of (\ref{eq:F}) now gives
\begin{equation}\label{eq:smoothing-bound}
 |\partial_\alpha^p\partial_\beta^qF(\alpha,\beta)|
 \leq\frac\pi2 S_{pq}(r),\qquad |\alpha|\leq r,\quad |\beta|\leq1.
\end{equation}
The differentiated series converges absolutely and uniformly on this
set, by (\ref{eq:Spq}); terms with insufficient indices vanish. Taking
$r=(1+\rho_0)/2$ and using $|P(z)|\leq\rho_0$ on $|z|\leq1$ proves
that every mixed derivative needed through order four is continuous
after composition on the closed disk. The chain rule consequently gives
a continuous boundary value for $D^4H$, with uniform convergence of
$D^4H(se^{i\theta})$ as $s\uparrow1$. This includes $z=1$ and $z=-1$.
This qualitative argument requires no numerical evaluation at the
parameter $(1+\rho_0)/2$.

We next give an explicit finite calculation of $S_{pq}(r)$ for
$1\leq p+q\leq4$. Put $T=\sqrt r\,w+L(x)$ and
$\psi(T)=\erf(T/\sqrt{2v})$. For $d\geq1$,
\[
 \psi^{(d)}(T)=\sqrt{2/\pi}\,v^{-1/2}e^{-T^2/(2v)}U_d(T),
 \qquad U_d(T)=(-1)^{d-1}v^{-(d-1)/2}\He_{d-1}(T/\sqrt v).
\]
In particular,
\[
 U_1=1,\quad U_2=-T/v,\quad U_3=T^2/v^2-1/v,
 \quad U_4=-T^3/v^3+3T/v^2.
\]
The chain rule gives
\begin{equation}\label{eq:R}
 \partial_w^p\partial_x^qG_r
 =r^{p/2}\sqrt{2/\pi}\,v^{-1/2}e^{-T^2/(2v)}R_{pq}(x,T),
 \quad R_{pq}=\sum_k B_{qk}U_{p+k}.
\end{equation}
Here $B_{00}=1$; for $q>0$ the nonzero Bell polynomials are
\begin{align*}
 B_{11}&=L',\\
 B_{21}&=L'',&B_{22}&=(L')^2,\\
 B_{31}&=L''',&B_{32}&=3L'L'',&B_{33}&=(L')^3,\\
 B_{41}&=L'''',&B_{42}&=4L'L'''+3(L'')^2,&
 B_{43}&=6(L')^2L'',&B_{44}&=(L')^4.
\end{align*}
These identities follow by differentiating the previous row and
collecting derivatives of $\psi$; no asymptotic chain rule is involved.
Only $p+k\geq1$ occurs when $p+q\geq1$.

For any polynomial $Q_0$ and fixed real $x$, completing the square in the
$w$ integral proves
\begin{equation}\label{eq:tilt}
 \E_w[e^{-T^2/v}Q_0(T)]
 =\sqrt{\frac v{1+r}}e^{-L(x)^2/(1+r)}\E[Q_0(Y)],
 \quad Y\sim N\left(\frac{vL(x)}{1+r},\frac{rv}{1+r}\right).
\end{equation}
Indeed the original density has variance $r$ and mean $L(x)$ in the
$T$ coordinate; adding $T^2/v$ to its exponent gives the stated mean,
variance, and normalizing factor. The moments of $Y$, with mean $\mu$
and variance $\sigma^2$, are exactly
\begin{equation}\label{eq:tilted-moments}
 \E Y^d=\sum_{\ell=0}^{\lfloor d/2\rfloor}
 \frac{d!}{2^\ell\ell!(d-2\ell)!}\sigma^{2\ell}\mu^{d-2\ell},
\end{equation}
as follows by expanding $\exp(\mu t+\sigma^2t^2/2)$.
With the real polynomial
$Q_{pq}(x)=\E[R_{pq}(x,Y)^2]$, equations
(\ref{eq:Spq})--(\ref{eq:tilted-moments}) give
\begin{equation}\label{eq:S-integral}
 S_{pq}(r)=\frac{2}{\pi\sqrt{1-r^2}}
       \int_\R Q_{pq}(x)e^{-L(x)^2/(1+r)}\dd\gamma(x).
\end{equation}
This polynomial is even, is nonnegative on the real line, and has degree
at most 118 in the cases used here. Its coefficients need not be
nonnegative. Evenness also follows directly from simultaneous sign
change of $w,x$ in the squared derivative in (\ref{eq:R}).

There is a second way to construct exactly the same polynomial, used in
the accompanying implementation. For polynomials in $T$ with polynomial
coefficients in $x$, let
\[
 \mathcal T R=\partial_T R-TR/v,\qquad
 \mathcal X R=\partial_xR+L'(x)\mathcal T R.
\]
If $p>0$, start from $1$, apply $\mathcal T$ a total of $p-1$ times,
then $\mathcal X$ a total of $q$ times. If $p=0$, start from $L'$ and
apply $\mathcal X$ a total of $q-1$ times. Differentiating
$e^{-T^2/(2v)}R$ verifies these rules directly. Squaring the result and
using (\ref{eq:tilted-moments}) gives $Q_{pq}$, providing a finite
construction independent of the Bell polynomial listing.

\subsection{The complete error for the smoothing integrals}
Every finite integral in (\ref{eq:S-integral}) is evaluated on $[0,16]$
as the sum of the 16 integrals on $[j,j+1]$, then doubled. The complex
integrand supplied to the interval integrator is the entire function
\begin{equation}\label{eq:entire-S}
 J(z)=\frac{2Q_{pq}(z)}{\pi\sqrt{1-r^2}\sqrt{2\pi}}
       \exp\left(-\frac{L(z)^2}{1+r}-\frac{z^2}{2}\right).
\end{equation}
In particular, the square in this formula is the analytic square, not a
complex absolute square. Enclosing polynomial coefficients enclose the
exact algebraic coefficients of $L$ and $Q_{pq}$.

The finite-interval calculation uses the certified complex integration
algorithm in Arb/FLINT~\cite{arb,integration,flintcalc}. Its adaptive
Gauss--Legendre rule encloses the integrand on a full complex ellipse;
it accepts a quadrature result only with a rigorous remainder. For an
ellipse of parameter $\rho>1$, maximum bound $M$, and $s$ nodes, the
bound used by that algorithm on $[-1,1]$ is
\begin{equation}\label{eq:library-error}
 \frac{64M}{15(\rho-1)\rho^{2s-1}}.
\end{equation}
It is the analytic Gauss--Legendre remainder bound implemented in the
cited algorithm; scaling by the interval half-width gives the bound on
a subinterval. A fully elementary, more conservative remainder
$8M\rho^{-2s}/(1-\rho^{-1})$ follows from the Chebyshev proof above and
also explains why such enclosing quadratures converge here. Subdivision
adds enclosing integrals and their errors, rather than extrapolating
point samples. Direct integration enclosures on small intervals are
also permitted by the algorithm. All calls use 384-bit arithmetic,
relative tolerance $10^{-18}$, absolute tolerance $10^{-20}$, and an
evaluation limit of $10^6$; these tolerances request accuracy and are
not themselves claimed as error bounds. Only the returned finite
enclosures are used. Entire analyticity of (\ref{eq:entire-S}) validates
every analyticity query and every ellipse bound required by the
integrator. The mathematical dependence here is the proved enclosing
quadrature algorithm and the inclusion property of ball arithmetic,
not an uncertified numerical integration routine.

The infinite remainder is bounded separately. Write
$L(x)=\sum_{k=0}^9\lambda_k x^k$. The coefficient enclosures prove
\begin{equation}\label{eq:Ltail}
 \sum_{k<9}|\lambda_k|16^{k-9}<\frac{|\lambda_9|}{2},
 \qquad c:=10^{-6}<\frac{|\lambda_9|}{2}.
\end{equation}
Hence $|L(x)|\geq c|x|^9$ for $|x|\geq16$. If
$m=\deg Q_{pq}$, the further checked inequality
\[
 \frac{18c^2 16^{18}}{1+r}>m
\]
shows that $x^ke^{-c^2x^{18}/(1+r)}$ decreases for $x\geq16$ and
every $0\leq k\leq m$. Therefore the entire two-sided omitted integral
in (\ref{eq:S-integral}) is at most
\begin{equation}\label{eq:S-tail}
 \frac{2}{\pi\sqrt{1-r^2}}
   \left(\sum_{k=0}^m|[x^k]Q_{pq}|16^k\right)
       e^{-c^2 16^{18}/(1+r)}<2^{-100}.
\end{equation}
In deriving this bound we used only that the Gaussian measure of
$\{|x|>16\}$ is at most one, so both tails are included. A symmetric
interval of this radius is added to twice the finite real integral.
All 266 required integral records contain the finite integral enclosure,
this tail bound, and the positive margins in (\ref{eq:Ltail}) and the
monotonicity test. These records are computed from an empty integral
cache.

\subsection{Mixed derivative bounds from complex Gaussian densities}
For $z\in\C$ with $|\Re z|<1$, the analytic continuation of a correlated
Gaussian density, after the real orthogonal change of variables
$U=(x+y)/\sqrt2$, $V=(x-y)/\sqrt2$, is
\[
 k_z(U,V)=\frac{1}{2\pi\sqrt{1+z}\sqrt{1-z}}
       \exp\left(-\frac{U^2}{2(1+z)}-\frac{V^2}{2(1-z)}\right).
\]
Each square root is its analytic branch in the right half-plane.
On compact subsets of the strip the real parts of the inverse variances
are positive and uniformly bounded below. Thus this density and every
fixed parameter derivative have an integrable Gaussian majorant times
a polynomial. Differentiation of its integral against bounded real
functions is justified by dominated convergence and defines a
holomorphic function.

Let $d_0(z)=\int_{\R^2}|k_z|$. Completing two real Gaussian integrals
gives
\begin{equation}\label{eq:density-mass}
 d_0(z)=\sqrt{\frac{|1+z|\,|1-z|}{1-(\Re z)^2}},\quad
 v_U=\frac{|1+z|^2}{1+\Re z},\quad
 v_V=\frac{|1-z|^2}{1-\Re z}.
\end{equation}
The probability density $|k_z|/d_0(z)$ is that of independent centered
real Gaussians with variances $v_U,v_V$. To compute derivatives write
$\partial_z^p k_z=k_z T_p(U,V;z)$. One explicit formula is
\begin{align}
 R_j(v,u)&=\sum_{\ell=0}^j
 \frac{(2j)!(-1)^{j-\ell}}{2^{2j-\ell}(j-\ell)!(2\ell)!}
                       v^{-j-\ell}u^{2\ell},\nonumber\\
 T_p&=\sum_{j=0}^p\binom pj(-1)^{p-j}
                          R_j(1+z,U)R_{p-j}(1-z,V).\label{eq:density-poly}
\end{align}
For a one-dimensional Gaussian density, the heat equation
$\partial_v=\tfrac12\partial_u^2$ and the Hermite formula for its
spatial derivatives prove $R_j$. The product rule and
$\partial_z(1-z)=-1$ then prove (\ref{eq:density-poly}).
Cauchy--Schwarz under the probability law just described gives
\begin{equation}\label{eq:density-bound}
 \norm{\partial_z^pk_z}_{L^1(\R^2)}
 \leq d_p(z):=d_0(z)\sqrt{\E|T_p(U,V;z)|^2}.
\end{equation}
All expectations here are finite polynomial Gaussian moments, without
any integral truncation.

For an alternative exact construction, the $k$th logarithmic derivative
$\ell_k=\partial_z^k\log k_z$ has constant, $U^2$, and $V^2$
coefficients respectively
\begin{align*}
 &\frac{(-1)^k(k-1)!}{2(1+z)^k}
           +\frac{(k-1)!}{2(1-z)^k},\\
 &\frac{(-1)^{k+1}k!}{2(1+z)^{k+1}},\qquad
           -\frac{k!}{2(1-z)^{k+1}}.
\end{align*}
Starting with $T_0=1$, the finite recurrence
\[
 T_m=\sum_{k=1}^m\binom{m-1}{k-1}\ell_kT_{m-k}
\]
is obtained by differentiating an exponential generating function for
$k_{z+t}/k_z$. It is the construction used in the independent
implementation. To evaluate its squared norm without a cancellation
assumption, rescale $U=\sqrt{v_U}\,u$, $V=\sqrt{v_V}\,v$, and use
\[
 u^{2k}=\sum_{i=0}^k
 \frac{(2k)!}{2^{k-i}(k-i)!(2i)!}\He_{2i}(u).
\]
This identity follows by comparing coefficients in the Hermite
generating function multiplied by $e^{t^2/2}$. If the resulting
polynomial is $\sum_{i,j}c_{ij}\He_{2i}(u)\He_{2j}(v)$, then
\[
 \E|T_p|^2=\sum_{i,j}|c_{ij}|^2(2i)!(2j)!.
\]
Thus all complex conjugations occur only in a finite nonnegative square
norm, not in the analytic density or its derivative.

Using independent copies of the correlated density for the $W$ and $X$
coordinates, boundedness $|fg|\leq1$ gives
\begin{equation}\label{eq:mixed-density}
 |\partial_\alpha^p\partial_\beta^qF(\alpha,\beta)|
 \leq\frac\pi2d_p(\alpha)d_q(\beta),
 \qquad |\Re\alpha|,|\Re\beta|<1.
\end{equation}
To identify this analytic function with (\ref{eq:F}), first note their
equality on the real open square of correlations. Fixing a real second
coordinate and using the one-variable identity theorem, then fixing a
complex first coordinate and repeating, proves equality in the open
bidisk. At boundary points within the strip, the derivatives agree with
the continuous limits established by (\ref{eq:Spq}). This proves that
(\ref{eq:mixed-density}) bounds the same derivatives used by the
coefficient argument.

\subsection{Angular chain rule and complete interval coverage}
For $s\geq0$ write
\[
 D^sH(z)=\sum_{p,q}c^{(s)}_{pq}(z)
               (\partial_\alpha^p\partial_\beta^qF)(P(z),z).
\]
These are exact rational polynomial coefficients. For example they are
given by the finite formula
\begin{equation}\label{eq:angular-jets}
 c^{(s)}_{pq}(z)=\frac{s!}{p!q!}[u^s]
             (P(ze^u)-P(z))^p(ze^u-z)^q.
\end{equation}
Taylor's formula for $F(P(ze^u),ze^u)$ proves this identity in the open
disk. It can also be computed recursively, with out-of-range indices
zero, from
\begin{equation}\label{eq:angular-rec}
 c^{(0)}_{00}=1,\qquad
 c^{(s+1)}_{pq}=D c^{(s)}_{pq}+(DP)c^{(s)}_{p-1,q}+z c^{(s)}_{p,q-1}.
\end{equation}
For $s=4$ the nonzero indices are precisely
$p,q\geq0$, $1\leq p+q\leq4$, a set of 14 elements. All these
polynomials are computed over $\mathbb Q$, so the finite chain rule
contains no numerical differentiation error.

Divide the first quadrant into the 512 closed intervals
\[
 I_j=[j\pi/1024,(j+1)\pi/1024],\qquad 0\leq j<512.
\]
Interval sine and cosine of $kI_j$ enclose $z^k=e^{ik\theta}$ for
every $\theta\in I_j$. This yields enclosures for $P(z)$ and all
$c^{(4)}_{pq}(z)$ on each whole interval. A smoothing bound is available
whenever the enclosure proves $|P(z)|<r$ for one of
$r=78/100,79/100,\ldots,97/100$. A density bound is available when
both $|\Re P(z)|<1$ and $|\Re z|<1$ are proved by the same interval
enclosures. Every panel has at least one available estimate. Where
both are available, their smaller upper bound may be used separately
for each $(p,q)$. The intervals touching $z=1$ are covered by smoothing;
the density estimate is not asserted at its singular endpoint.

Let $U_{j,pq}$ be such an upper bound for
$(2/\pi)|\partial_\alpha^p\partial_\beta^qF|$, and let $C_{j,pq}$
bound $|c^{(4)}_{pq}|$. Then
\begin{equation}\label{eq:panel-majorant}
 |D^4H(e^{i\theta})|\leq M_j:=\frac\pi2
             \sum_{1\leq p+q\leq4}C_{j,pq}U_{j,pq}
             \quad(\theta\in I_j).
\end{equation}
The circle certificate records all $512\cdot14=7168$ terms, each
chosen method, and its domain inequalities. Exactly 266 smoothing
integrals are required: all 14 indices at each of the 19 parameters
$79/100,\ldots,97/100$. These are the integrals bounded in
(\ref{eq:S-integral})--(\ref{eq:S-tail}).

Real coefficients and oddness imply
$D^4H(\bar z)=\overline{D^4H(z)}$ and $D^4H(-z)=-D^4H(z)$.
The four quadrants therefore give equal contributions to its squared
norm. Round each $M_j$ upwards to $\widehat M_j=k_j/10^8$ with
integer $k_j$. The complete list of these 512 rational upper bounds is
included in the certificate. Summing their exact rational squares gives
\begin{equation}\label{eq:rational-square}
 \frac1{2\pi}\int_0^{2\pi}|D^4H(e^{i\theta})|^2\dd\theta
 \leq\frac1{512}\sum_{j=0}^{511}\widehat M_j^2
 =Q:=\frac{1567474233421402382277578179}{2560000000000000000}.
\end{equation}
The explicit strict rational comparison
\begin{equation}\label{eq:Cmargin}
 C_*^2-Q=\frac{107039174409781821}{2560000000000000000}>0
\end{equation}
proves (\ref{eq:C4}). The bounds concern complete angular intervals,
not samples of the circle.

\section{The coefficient tail and the exact numerical conclusion}
For $0<s<1$, ordinary Fourier Parseval applied to the power series gives
\[
 \frac1{2\pi}\int_0^{2\pi}|D^4H(se^{i\theta})|^2\dd\theta
       =\sum_{n\geq0}n^8|b_n|^2s^{2n}.
\]
The continuous boundary extension proved above permits passage to the
limit on the left by uniform convergence. Monotone convergence of the
nonnegative series gives the limit on the right. Hence
\begin{equation}\label{eq:parseval-boundary}
 \sum_{n\geq0}n^8|b_n|^2
 =\frac1{2\pi}\int_0^{2\pi}|D^4H(e^{i\theta})|^2\dd\theta<C_*^2.
\end{equation}
For every positive odd $N$, Cauchy--Schwarz and the decreasing function
$(N+2x)^{-8}$ show that
\begin{align}
 \sum_{\substack{n>N\\n\ {\rm odd}}}|b_n|
 &\leq\left(\sum_{n>N}n^8|b_n|^2\right)^{1/2}
        \left(\sum_{k=1}^{\infty}(N+2k)^{-8}\right)^{1/2}\nonumber\\
 &<\frac{C_*}{\sqrt{14N^7}},\label{eq:tail-final}
\end{align}
because the latter sum is at most
$\int_0^\infty(N+2x)^{-8}\dd x=1/(14N^7)$.
For $N=501$ the first omitted coefficient is $b_{503}$.

Put
\[
 T_*=0.000002349547147305,\qquad \gamma_*=0.8818276493988.
\]
All these displayed terminating decimals denote exact rational numbers.
The certificate checks the following integer comparison:
\begin{equation}\label{eq:tail-rational}
 T_*^2(14\cdot501^7)-C_*^2
 =\frac{3104661541167915364988902768147}
        {20000000000000000000000000000000000}>0.
\end{equation}
Thus (\ref{eq:head-bounds}) and (\ref{eq:tail-final}) give
\begin{align}
 \Gamma
 &> B-h_*-T_*\nonumber\\
 &=0.88182764939885093252
   =\frac{22045691234971273313}{25000000000000000000}
   >\gamma_*,\label{eq:gamma-exact}\\
 B-h_*-T_*-\gamma_*
 &=\frac{1273313}{25000000000000000000}>0.\nonumber
\end{align}
In particular, the coefficient margin is positive with a fixed uniform
constant, independent of the dimensions or entries of $M$.

We give an elementary rational verification of the last reciprocal
comparison, so it does not depend on a decimal approximation to $\pi$.
Machin's identity is
\[
 \pi=16\arctan(1/5)-4\arctan(1/239).
\]
It follows, for example, by applying the tangent addition law to
$4\arctan(1/5)-\arctan(1/239)$ and observing that this angle lies in
$(0,\pi/2)$ and has tangent one. For $0<x<1$, let
$s_{64}(x)=\sum_{k=0}^{63}(-1)^kx^{2k+1}/(2k+1)$. The alternating
series remainder gives
\[
 s_{64}(x)<\arctan x<s_{64}(x)+x^{129}/129.
\]
Substitution at $1/5$ and $1/239$ proves, by rational arithmetic,
\[
 \pi<\pi_U:=3.141592653589793238462643383280.
\]
Finally,
\begin{equation}\label{eq:reciprocal}
 2\gamma_*\,(1.781296297373)-\pi_U
 =\frac{8045255950767709}{12500000000000000000000000000}>0.
\end{equation}
Lemma~\ref{lem:rounding}, (\ref{eq:gamma-exact}), and
(\ref{eq:reciprocal}) prove Theorem~\ref{thm:main}. Taking the supremum
over finite matrices is legitimate because the same constant
$\pi/(2\gamma_*)$ works for every matrix, and that constant is strictly
less than the displayed decimal.

The degree-501 calculation also proves the conservative bound
$K_G^{\R}<1.78130$ using only $C_4<27159$ in place of
(\ref{eq:C4}). Indeed, rational comparisons give
\[
 \frac{27159}{\sqrt{14\cdot501^7}}<0.000002578800327634,
 \quad B-h_*-0.000002578800327634>0.8818274201456,
\]
and $\pi/(2\cdot0.8818274201456)<1.781296760466<1.78130$.
The sharper stated bound uses specifically (\ref{eq:rational-square})
and (\ref{eq:Cmargin}).

\appendix
\section{The exact preprocessing polynomial}\label{app:polynomial}
The polynomial is $P(t)=10^{-18}\sum_{d=1,3,\ldots,81}a_dt^d$, with
the following integer coefficients. Every coefficient not listed is zero.
The second, independent preprocessing channel is exactly $B(t)=t$.
\begin{center}
\small
\begin{tabular}{rr@{\qquad}rr@{\qquad}rr}
\toprule
$d$ & $a_d$ & $d$ & $a_d$ & $d$ & $a_d$ \\
\midrule
1 & 890309324376476918 & 29 & -13910936207 & 57 & 832954101 \\
3 & -107057575215619705 & 31 & 5917793946 & 59 & 605857643 \\
5 & 2593131841827897 & 33 & 5403085315 & 61 & -74510125 \\
7 & 921687192873 & 35 & -4149053342 & 63 & -272828457 \\
9 & -18860525521356 & 37 & -1381031720 & 65 & 80076459 \\
11 & 9201650732787 & 39 & 2659882186 & 67 & 362652265 \\
13 & 181911301762 & 41 & -76041523 & 69 & 231024103 \\
15 & -1094634337148 & 43 & -2832277366 & 71 & -62280446 \\
17 & -42579971855 & 45 & -1408197678 & 73 & -164242940 \\
19 & 102427747255 & 47 & 885884647 & 75 & -53524188 \\
21 & -178837625035 & 49 & 947390899 & 77 & 66231724 \\
23 & -91902676358 & 51 & -335823134 & 79 & 73563652 \\
25 & 32692014381 & 53 & -760243970 & 81 & 12013495 \\
27 & 4095525653 & 55 & 71576160 &  &  \\
\bottomrule
\end{tabular}
\end{center}
The five rational coefficients of $L$ are specified completely in
(\ref{eq:L}), in the orthonormal probabilists\textquotesingle{} Hermite basis.
Thus neither a floating-point optimization output nor a choice of Hermite
normalization is left implicit in these data.

\section{Certificate contents and reproduction}
The accompanying directory \nolinkurl{reproducibility/} contains the exact
input, the numerical programs, and the complete finite enclosures. Its
\nolinkurl{immutable_manifest.json} specifies SHA-256 hashes for every
unchanged numerical input, program, and certificate file. The following
dependencies distinguish finite data from mathematical hypotheses.
\begin{enumerate}
\item The exact input is \nolinkurl{head/sources/candidate_rational.json}.
Its polynomial coefficients are printed in (\ref{eq:L}) and
Appendix~\ref{app:polynomial}. Its SHA-256 hash is
\begin{quote}\small\ttfamily
207b7ed03d6bc1fd44922863e361544bb87442d13801870f0a8bc6705e8a3040.
\end{quote}
Only the exact rational coefficient strings determine the polynomials;
diagnostic fields in the file are not assumptions.
\item The moment calculation uses \nolinkurl{head/run/gauss_nodes_112.json},
\nolinkurl{head/run/head501_config.json}, and all 32 files in
\nolinkurl{head/run/blocks/}. The merged file is
\nolinkurl{head/run/head_moments_501.json}. Its full index set is
$\mathcal I$; the complete quadrature and infinite remainders are
(\ref{eq:uniform-error}) and (\ref{eq:moment-tail}). The configuration
binds the input, source, node, and block parameters. Each block is bound
to that configuration and contains exactly its prescribed four panels.
\item The two complete composed coefficient files are
\nolinkurl{head/run/forward_head_501.json} and
\nolinkurl{head/evidence/independent_head.json}. Exact rational powers
and the index proof in (\ref{eq:head-formula}) establish that no
coefficient through degree 501 is missing. Moment arithmetic and the
rational-power assembly use 384 bits; the separate interval-series
assembly uses 160 bits. The exact node containment calculation uses
rational arithmetic and does not accept approximate root locations.
\item The norm calculation uses the files
\nolinkurl{independent_integrals.json} and
\nolinkurl{independent_circle.json} in \nolinkurl{norm/evidence/}.
These contain all 266 integrals and all 512 closed panels with their
14 terms, including the domain inequalities for each estimate.
Equations (\ref{eq:S-integral})--(\ref{eq:S-tail}) justify the integral
values and their complete errors; equations
(\ref{eq:density-mass})--(\ref{eq:panel-majorant}) justify every other
finite term. The filename and the auxiliary degree-251 tail field of
the norm engine do not limit its derivative bound: its integral and
circle calculations depend only on $L$, $P$, and derivative order four.
The degree-501 tail is calculated separately from (\ref{eq:tail-final}).
\item The file \nolinkurl{rational_certificate.json} contains every
integer $k_j$ in (\ref{eq:rational-square}), the exact rational head
endpoints, and all the rational inequalities
(\ref{eq:Cmargin}) and (\ref{eq:tail-rational})--(\ref{eq:reciprocal}).
It binds the coefficient, integral, and circle files by their hashes.
The rational caps are at least both the saved panel majorants and
independently re-summed products of all their saved term enclosures.
\end{enumerate}

The arithmetic implementation uses Python, python-flint 0.8.0, and
FLINT 3.3.1. Its mathematical requirements are inclusion-preserving
real and complex ball operations, correct exact rational arithmetic,
and the certified integral algorithm described above. Floating-point
Legendre initializers, approximate optimization data, stored success
flags, and prior integral caches are not premises. SciPy is optional:
if unavailable, stored node midpoints provide initial guesses, after
which all brackets, signs, and weights are proved again.

From the directory containing \nolinkurl{verify.py}, the command
\begin{quote}\small\ttfamily
python -B -u verify.py --output /tmp/real-upper-verification
\end{quote}
checks hashes and numerical bindings, reproduces the complete moment
merge and the uniform error, verifies the nodes, reruns both coefficient
assemblies, and recomputes all smoothing integrals and circle panels
from an empty integral cache. It then checks the exact rational
conclusion. Recomputed integral balls may have different endpoints
because of adaptive subdivision or arithmetic versions; each must be a
valid enclosure and must itself give $C_4<C_*$. In particular, the
printed value of $Q$ specifies the supplied rational certificate and
does not require bitwise equality of independently recomputed
quadrature results.

The separate program \nolinkurl{reproduce.py} initializes new nodes and
quadrature errors and recomputes every moment block, with one file per
block so that interruptions can be resumed. It then performs both
assemblies, a new integral and circle calculation, and the rational
comparison. It neither installs software nor acquires computational
resources. Its command-line options and complete reproduction commands
are in \nolinkurl{README.txt} and \nolinkurl{commands.txt}. The proof of the
infinite coefficient bound is the regularity and Parseval argument in
the main text; no finite file or numerical sample substitutes for that
argument.

\begin{thebibliography}{9}
\bibitem{krivine}
J.-L. Krivine,
\emph{Constantes de Grothendieck et fonctions de type positif sur les sph\`eres},
S\'eminaire Maurey--Schwartz 1977--1978, expos\'es 1 et 2, pp. 1--17.
\url{https://www.numdam.org/item/SAF_1977-1978____A1_0/}.

\bibitem{bmmn}
M. Braverman, K. Makarychev, Y. Makarychev, and A. Naor,
\emph{The Grothendieck constant is strictly smaller than Krivine's bound},
2011. \url{https://arxiv.org/abs/1103.6161}.

\bibitem{saha}
R. Saha, A. Li, A. Xue, S. Chaudhuri, A. Klivans, P. K. Kothari,
and R. Meka,
\emph{New Lower and Upper Bounds for the Grothendieck Constant},
arXiv:2608.11158, version 2.
\url{https://arxiv.org/html/2608.11158v2}.

\bibitem{li}
A. Li, R. Saha, A. Xue, S. Chaudhuri, A. Klivans, P. K. Kothari,
and R. Meka,
\emph{Long-Horizon AI Research for Grothendieck Constant:
A Case Study in Human--AI Mathematical Collaboration},
arXiv:2608.11195, version 3.
\url{https://arxiv.org/html/2608.11195v3}.

\bibitem{arb}
F. Johansson, \emph{Arb: Efficient Arbitrary-Precision Midpoint-Radius
Interval Arithmetic}, IEEE Transactions on Computers \textbf{66} (2017),
1281--1292. \href{https://doi.org/10.1109/TC.2017.2690633}
{doi:10.1109/TC.2017.2690633}.

\bibitem{integration}
F. Johansson, \emph{Numerical integration in arbitrary-precision ball
arithmetic}, in Mathematical Software---ICMS 2018, Lecture Notes in
Computer Science \textbf{10931}, Springer, 2018, pp. 255--263.
\url{https://arxiv.org/abs/1802.07942}.

\bibitem{flintcalc}
The FLINT developers, \emph{acb\_calc: calculus with complex functions},
FLINT documentation, integration and Gauss--Legendre remainder routines.
\url{https://flintlib.org/doc/acb_calc.html}.
\end{thebibliography}

\end{document}
